
\begin{abstract}
Input-conditioned low-rank modulation combines compact persistent factors with
a generated diagonal core. Its representational, storage, and execution effects
must be distinguished to establish an efficiency benefit. We study this
separation for gated Transformer feed-forward networks (FFNs). We derive an
output-subspace error floor that no controller can overcome while the final
factor remains fixed, an activation-weighted constant-core control, and a
matched-schedule bound on removable descriptor transfers. These results give
falsifiable tests of when dynamic modulation and descriptor fusion are useful.
We then reanalyze archived accelerator records and perform fresh software
verification. In four recorded training-seed summaries for a Qwen2.5-1.5B FFN,
the generated model improves relative reconstruction MSE only marginally at
matched rank and is 1.95\% worse at the closest reported parameter budget.
Original execution reports and checkpoints for that family are unavailable,
so these are auditable archive calculations rather than independently
reproduced quality measurements. A same-run archived kernel summary reports
0.8234 versus 0.8509 ms for fused and materialized execution at 128 tokens,
but persistent low rank takes 0.2719 ms. A later multicore variant fails its
recorded repeated-call stability gate. Fresh CPU tests of the actual causal
FFN reduce forward-retained nonparameter tensor storage by $13\times$ under
checkpointing, with identical outputs and gradients in nine small synthetic
fixtures. The evidence supports precise mechanism and limitation statements;
it does not yet establish a full-model quality--latency advantage.
\end{abstract}

\section{Introduction}
Transformer feed-forward networks contain large parameter matrices and several
intermediate activations. Reducing one of these objects can change arithmetic,
parameter capacity, memory transfers, and kernel launch overhead simultaneously.
A speedup observed after a low-rank replacement therefore does not identify
which change caused it. This distinction matters for conditional weight
generation: a compact controller may add expressivity, while computing and
materializing its outputs can cost more than the transfers that fusion removes.

We examine \mwg, a causal, token-conditioned diagonal modulation of persistent
low-rank FFN factors. For a row vector $x$, one projection is
\begin{equation}
 p(x)=\bigl((xU)\odot z(x)\bigr)V,
 \quad U\in\R^{d_{\rm in}\times r},\quad V\in\R^{r\times d_{\rm out}}.
 \label{eq:projection}
\end{equation}
The algebra is closely related to feature modulation and hypernetworks.
The systems question is whether the generated descriptor $z$ should be stored
as a global tensor or consumed immediately where it is formed. We call the
latter a \emph{descriptor no-store contract}. The factors, controller parameters,
and output activations retain their own lifetimes; the contract concerns only
the named descriptor and a declared implementation boundary.

Three contributions organize this study. First, we derive representation and
control results that separate changing rank from changing conditional
coefficients: an output-subspace obstruction and an optimal constant-core
projection control. Second, we give a memory-interface-specific accounting
that separates parameter storage, generated cardinality, scheduled global
transfers, and physical DRAM traffic. Third, we evaluate the resulting claims
against a reproducible archive audit and new CPU software checks. Negative
comparisons and failed validation gates are retained because they determine
which deployment conclusions the evidence can support.

The most useful empirical conclusion is conditional. Generated modulation can
approximately match a trained same-rank reconstruction baseline in the available
summaries, but the strongest same-budget comparison favors the persistent
baseline. Eliminating an explicit descriptor tensor is feasible at the source
level; it does not imply that the complete FFN runs faster. Both conclusions
are consistent with the analysis and motivate a smaller set of decisive
experiments than a broad collection of disconnected speed and quality claims.

\section{Relationship to Prior Work}
\paragraph{Conditional parameters and feature modulation.}
HyperNetworks generate another network's parameters \citep{ha2017hypernetworks}.
FiLM applies conditional feature-wise affine transformations
\citep{perez2018film}; Eq.~\eqref{eq:projection} is multiplicative modulation in
a rank-space representation. Learned static activation scales also appear in
$(\mathrm{IA})^3$ \citep{liu2022ia3}. A static diagonal scale can be absorbed
into a factor, whereas a token-conditioned scale generally cannot. LoRA
\citep{hu2022lora} is a low-rank update to a retained base matrix, while the
studied FFN replaces each dense projection by factors. These distinctions
preclude treating the present algebra as an unprecedented form of generation.

The recent Ouroboros preprint uses a hidden-state controller to modulate
SVD-initialized low-rank bases within recursive Transformers
\citep{jaber2026ouroboros}. It is especially close algebraic prior art.
Our scope is the lifetime and controlled assessment of FFN descriptors,
including the limits imposed by a fixed output subspace. We make no claim of
priority over input-conditioned diagonal modulation.

\paragraph{I/O-aware execution and recomputation.}
FlashAttention analyzes and reduces transfers between memory levels
\citep{dao2022flashattention}. FlashMHF combines dynamically weighted FFN
subnetworks with online SRAM computation \citep{zhang2025flashmhf}; it is a
relevant future systems and model-quality comparator. Activation checkpointing
trades retained tensors for recomputation \citep{chen2016checkpoint}, and
Checkmate optimizes rematerialization schedules \citep{jain2020checkmate}.
Our CPU checkpoint experiment applies this established technique to the actual
MWG module. It verifies compatibility and instrumentation rather than proposing
a new checkpointing algorithm.

\paragraph{Compression and fair comparison.}
Quantization methods such as GPTQ and AWQ reduce the representation cost of
weights \citep{frantar2023gptq,lin2024awq}. They remain important deployment
baselines. Rank reduction, quantization, modulation, and fusion should be
compared at matched task quality and reported budgets. The archived experiments
here do not support numerical superiority over those methods. The output
rank bound uses the classical Eckart--Young--Mirsky approximation result
\citep{eckart1936approximation,mirsky1960norm}.

\begin{table}[t]
\centering\small
\caption{Which mechanism a controlled contrast can identify.}
\label{tab:contrast-design}
\begin{tabular}{p{.22\linewidth}p{.34\linewidth}p{.34\linewidth}}
\toprule
Question & Matched contrast & Required invariant \\
\midrule
Does rank reduction help? & Dense versus trained persistent factors & Task metric, data and optimization budget \\
Does conditioning help? & Dynamic versus constant-core/static control & Parameter or compute budget; causal evaluation \\
Does descriptor fusion help? & Same dynamic model, materialized versus fused & Checkpoint, inputs, output tolerance and timing boundary \\
Does recomputation help? & Eager versus checkpointed same module & Parameters, input, loss and all gradients \\
\bottomrule
\end{tabular}
\end{table}

\section{Model and Execution Contract}
\subsection{Causal gated FFN}
Let $d$ be the model width and $m$ the FFN intermediate width. The dense
reference, using row-vector conventions, is
\begin{equation}
 f(x)=\bigl[\operatorname{SiLU}(xW_g)\odot(xW_u)\bigr]W_d,
 \quad W_g,W_u\in\R^{d\times m},\quad W_d\in\R^{m\times d}.
\end{equation}
In the factorized model, each $W_p$ is replaced by Eq.~\eqref{eq:projection}.
For token state $x_{b,t}$ the controller produces all three rank vectors:
\begin{align}
 h_{b,t}&=\operatorname{SiLU}(x_{b,t}A+a_0),\\
 (z_u,z_g,z_d)_{b,t}&=\mathbf 1+\alpha\tanh(h_{b,t}C+c_0),\label{eq:controller}\\
 \widetilde f(x)&=p_d\bigl(\operatorname{SiLU}(p_g(x))\odot p_u(x);z_d(x)\bigr).
 \label{eq:ffn}
\end{align}
Here $A\in\R^{d\times h}$ and $C\in\R^{h\times3r}$; all factors and
controller parameters are trainable persistent state. The down-projection
core is generated from the original pre-FFN state $x$, while its projection
input is the gated intermediate. The controller has
\begin{equation}
 P_G=(d+1)h+(h+1)3r\label{eq:controller-count}
\end{equation}
parameters, including biases. For equal fp32 factor/controller precision, total
parameter bytes are $4[3r(d+m)+P_G]$. Optimizer states and saved activations
are separate quantities.

The implemented causal mode uses a distinct core for each token ($q=1$).
It reads no future hidden states and pools no companion examples. If a future
implementation shares a core across $q>1$ tokens, both group boundaries and
its context must be determined from information available by the first token's
pre-FFN state. Pooling the complete group's teacher-forced states would violate
this condition. A batch of $B$ examples of length $S$ has
\begin{equation}
 n=B\lceil S/q\rceil,\qquad 1\le q\le S,\label{eq:groups}
\end{equation}
groups under batch-isolated grouping. Cached decoding has $S=q=1$ and $n=B$;
we do not label a 128-token call as single-token autoregressive decoding.

\subsection{The actual fusion boundary}
The custom Ascend operator receives rank activations $a=xU$, controller hidden
states $h$, and the controller head's parameters. It forms the head affine
result, applies the bounded transform in Eq.~\eqref{eq:controller}, and
multiplies $a$ by the resulting scale. The descriptor is not returned as an
explicit global tensor. The operator \emph{does} write $a\odot z$ to global
memory, and the $U$ and $V$ GEMMs run outside it. Consequently, neither the
entire projection nor the entire FFN is scratch-resident.

\begin{figure}[t]
\centering
\begin{tikzpicture}[>=Latex,node distance=10mm,
 box/.style={draw,rounded corners,align=center,font=\small,minimum height=10mm}]
\node[box] (input) {$x$\\pre-FFN state};
\node[box,right=of input] (a) {$a=xU$\\GEMM};
\node[box,right=of a,fill=blue!5] (fuse) {$z=1+\alpha\tanh(hC+c_0)$\\$a'=a\odot z$};
\node[box,right=of fuse] (out) {$y=a'V$\\GEMM};
\node[box,below=14mm of a] (h) {$h=\operatorname{SiLU}(xA+a_0)$};
\draw[->] (input)--(a);
\draw[->] (a)--node[above,font=\scriptsize]{global $a$}(fuse);
\draw[->] (fuse)--node[above,font=\scriptsize]{global $a'$}(out);
\draw[->] (input.south)|-(h.west);
\draw[->] (h.east)-|node[pos=.2,below,font=\scriptsize]{global $h$}(fuse.south);
\node[above=3mm of fuse,font=\small] {custom operator};
\end{tikzpicture}
\caption{Implemented boundary for one projection. Only the descriptor $z$ is
consumed without an explicit global descriptor tensor. Source-level absence
of a descriptor store does not prove absence of compiler or runtime spills.}
\label{fig:boundary}
\end{figure}

We distinguish three verification levels. A source audit checks that $z$ is not
returned, aliased, or explicitly stored. A compiled-kernel inspection checks
workspace and spill behavior for the exact binary. A device-counter comparison
measures transfers under a declared workload and cache state. The first does
not imply the other two. Binary hashes bind inspected code to an artifact;
they do not prove that a reported remote run actually used it.

During training, autograd may retain controller intermediates and descriptors.
The inference no-store contract therefore does not automatically extend to
backward. Deterministic rematerialization is one possible training strategy;
Section~\ref{sec:cpu} tests software checkpointing, not the custom kernel's
backward implementation.

\section{Representation Limits and a Stronger Control}
\subsection{Fixed output-subspace obstruction}
A conditional core changes coefficients in fixed rank directions. In one
projection, the effective matrices belong to the linear span of the $r$
outer products formed by columns of $U$ and rows of $V$; the span has dimension
at most $r$. The input-output mapping may still be nonlinear because $z$
depends on $x$. The final output restriction is stronger and independent of
the controller's complexity.

\begin{theorem}[Output-subspace obstruction]\label{thm:subspace}
Let $V_d\in\R^{r\times d}$ be the fixed final factor in Eq.~\eqref{eq:ffn},
let $\mathcal S$ be its row space, and let $P_{\mathcal S}$ be the orthogonal
projector onto $\mathcal S$. Stack $N$ target FFN outputs into $Y\in\R^{N\times d}$
and corresponding model outputs into $\widehat Y$. Then, for any controller,
\begin{equation}
 \norm{Y-\widehat Y}_F^2
 =\norm{Y(I-P_{\mathcal S})}_F^2
  +\norm{YP_{\mathcal S}-\widehat Y}_F^2
 \ge\norm{Y(I-P_{\mathcal S})}_F^2.\label{eq:subspace-floor}
\end{equation}
If the final factor is optimized subject to rank at most $r$, every candidate
still satisfies
\begin{equation}
 \norm{Y-\widehat Y}_F^2\ge\sum_{j>r}\sigma_j(Y)^2.\label{eq:rank-floor}
\end{equation}
\end{theorem}
\begin{proof}
Every output row is a coefficient row multiplied by $V_d$, so
$\widehat YP_{\mathcal S}=\widehat Y$. The two terms in the decomposition
$Y-\widehat Y=Y(I-P_{\mathcal S})+(YP_{\mathcal S}-\widehat Y)$ have orthogonal
row components, giving the Pythagorean identity. Also
$\rank(\widehat Y)\le r$; the best rank-$r$ approximation of $Y$ has squared
Frobenius error equal to its singular-value tail
\citep{eckart1936approximation,mirsky1960norm}.
\end{proof}

The theorem is a structural consequence of classical low-rank approximation,
not a new general approximation theorem. It explains a concrete limitation:
context-dependent coefficients cannot recover target components orthogonal to
the final persistent factor. The lower bound is evaluated on \emph{uncentered}
target outputs. For residual-block outputs, subtract the unchanged skip
connection before applying it. An arbitrary coefficient matrix can attain the
best projection onto a fixed subspace, but the restricted FFN/controller need
not attain it. The theorem does not say that controllers are useless within
that subspace, nor that all layers of a residual Transformer share one subspace.

A useful future diagnostic reports three quantities on fixed evaluation
activations: actual reconstruction error, the fixed-factor floor in
Eq.~\eqref{eq:subspace-floor}, and the best rank-$r$ output floor in
Eq.~\eqref{eq:rank-floor}. Their gaps distinguish coefficient approximation,
subspace selection, and unavoidable rank loss. These real-model floors cannot
be recovered from scalar MSE summaries alone and are not reported as measured
here.

\subsection{An activation-weighted constant-core baseline}
A shuffled context control changes the input to the controller and may change
the optimization process. It does not by itself show whether a learned dynamic
core can be replaced by an efficient constant. A sharper projection-level
control fits a constant using training/calibration activations only.

\begin{proposition}[Optimal frozen-projection constant]\label{prop:constant}
For fixed row-vector projection inputs $a=xU$, fixed $V$, and frozen dynamic
outputs $z(x)$, write $z(x)$ and $c$ as column vectors in this proposition.
Define $Q(x)=\diag(a)VV^\top\diag(a)$,
$H=\E[Q(x)]$, and $b=\E[Q(x)z(x)]$. Assume the needed second moments are finite.
The constants minimizing
\begin{equation}
 R(c)=\E\bigl[\norm{(a\odot(z(x)-c)^\top)V}_2^2\bigr]\label{eq:constant-risk}
\end{equation}
solve $Hc=b$. A minimum-norm solution is $c^*=H^\dagger b$; all other
solutions add a vector in $\ker(H)$.
\end{proposition}
\begin{proof}
$Q$ is positive semidefinite and
$R(c)=\E[z^\top Qz]-2c^\top b+c^\top Hc$.
If $v\in\ker(H)$ then $\E[v^\top Qv]=0$, so $Qv=0$ almost surely and
$v^\top b=0$. Thus $b$ belongs to the range of $H$, the normal equations
are consistent, and completing the quadratic gives the stated minimizers.
\end{proof}

The empirical version replaces expectations by calibration averages; its
constant folds into $U\diag(c^*)$. An ordinary mean of $z$ ignores the
activation and output metric $Q$. This is a diagnostic application of weighted
least squares, with no claim of novelty for the normal equations. Its optimum
is calibration-local, with no automatic held-out guarantee. In the three-stage
FFN, replacing earlier cores changes later inputs, so separately fitted
projection controls require a new full-FFN and full-model evaluation. No
constant-core quality improvement is inferred from the present archive.

\subsection{Sensitivity and invariances}
Let $A=\phi(XW_g)$, $B=XW_u$, and tildes denote approximations fixed within
one group. If $\phi$ is $L_\phi$-Lipschitz,
$\norm{W_p-\widetilde W_p}_F\le\delta_p$,
$\norm{B}_{\max}\le M_u$, and $\norm{\widetilde A}_{\max}\le M_g$, then
\begin{align}
 \norm{Y-\widetilde Y}_F
 &\le \norm{W_d}_2\norm{X}_2(M_uL_\phi\delta_g+M_g\delta_u)
   +\norm{\widetilde A\odot\widetilde B}_2\delta_d.\label{eq:sensitivity}
\end{align}
To verify it, use $C-\widetilde C=(A-\widetilde A)\odot B+
\widetilde A\odot(B-\widetilde B)$ and
$Y-\widetilde Y=(C-\widetilde C)W_d+
\widetilde C(W_d-\widetilde W_d)$, then apply submultiplicativity.
Token-conditioned matrices permit pointwise bounds followed by aggregation;
a fixed-matrix SVD optimum alone does not establish conditional learnability.

Two paired rank scales are redundant since $s\odot s'$ is a single scale.
Removing them does not make the complete factorization identifiable. For
invertible diagonal $D$, $U\mapsto UD$ and $V\mapsto D^{-1}V$ preserve
Eq.~\eqref{eq:projection} because diagonal matrices commute. All storage
comparisons count actual parameters rather than counting equivalent
parameterizations as distinct capabilities.

\section{Memory Accounting at a Declared Interface}
\subsection{Four quantities that must remain separate}
Dense parameter cardinality is $P_D=3dm$, and persistent factor cardinality
is $P_F=3r(d+m)$. A call generates $D_z=3nr$ diagonal scalars. These are
counts of objects; adding them does not produce a lower bound on DRAM traffic.
For example, a warm cache can suppress DRAM reads of persistent factors, while
multiple tiles can load the same factor repeatedly.

Let $b_w,b_G,b_z$ be bytes per factor, controller parameter, and descriptor
scalar. Parameter storage is $b_wP_F+b_GP_G$. Let $L_{\rm common}$ denote
transfers across one declared global-memory interface that are shared by two
implementations. It includes persistent-factor loads, controller loads,
source activations, scaled rank activations, and outputs according to the
actual schedule. It is not interchangeable with parameter storage. Physical
DRAM bytes require device counters and depend on caching and writeback.

\begin{proposition}[Removable final-descriptor transfers]\label{prop:bytes}
Suppose two schedules are identical except that one writes and reloads each
of the $3nr$ final descriptor scalars exactly once while the other consumes
those values locally. With $C_z=3nrb_z$,
\begin{equation}
 L_{\rm mat}=L_{\rm common}+2C_z,\qquad
 L_{\rm fused}=L_{\rm common}.\label{eq:bytes}
\end{equation}
If the common schedule includes one full cold read of every factor,
$L_{\rm common}\ge3r(d+m)b_w>0$, and therefore
\begin{equation}
 \frac{L_{\rm mat}}{L_{\rm fused}}
 \le1+\frac{2n b_z}{(d+m)b_w}.\label{eq:byte-ceiling}
\end{equation}
\end{proposition}
\begin{proof}
The only differing transfers are a store and load of $C_z$ bytes each.
Divide $L_{\rm common}+2C_z$ by $L_{\rm common}$ and substitute its lower bound.
\end{proof}

This is a scheduled-byte statement, not a physical-HBM or latency bound.
Different tiling, cache reuse, launch counts, controller-logit temporaries, casts,
or arithmetic violate the matched-schedule premise. In particular, an eager
materialized implementation may store more intermediates than the final
$z$; a fused implementation can remove those too. The bound isolates only the
final-descriptor round trip. It prevents attributing all improvements from
an unrelated multi-operation rewrite to one missing tensor.

\begin{figure}[t]
\centering
\includegraphics[width=.86\linewidth]{\paperroot/figs/journal_bytes.pdf}
\caption{Analytic upper ratio in Eq.~\eqref{eq:byte-ceiling}, assuming equal
factor/descriptor widths and one cold full-factor read. This is not measured
speedup. For $d+m=18432$, the ceilings are $1.01389$ at $n=128$ and
$1.44444$ at $n=4096$. Rank cancels in this isolated model.}
\label{fig:bytes}
\end{figure}

The corrected regime-map implementation reports scratch cardinality separately.
Its partial no-store global ledger is
$e_Fb_wP_F+e_Gb_GP_G+c_xBSd\,b_x$; materialization adds $2C_z$.
Common output and other activation transfers are not included, so these partial
totals must not be presented as whole-FFN traffic. The implemented ledger uses
$b_z=b_w$; Eq.~\eqref{eq:byte-ceiling} allows distinct widths. Version~2
explicitly changes the previous mixed-interface totals.

\subsection{Full-factor generation and reuse}
Generating all factor pairs per group emits $3nr(d+m)$ scalars instead of
$3nr$. Relative to the dense parameter cardinality, the one-copy ratio is
$n\gamma$, where $\gamma=r(d+m)/(dm)$. For positive integer dimensions and
$0<\gamma<1$, let $n_{\max}=\lceil1/\gamma\rceil-1$ and
$g_{\max}=\lfloor n_{\max}/B\rfloor$. Strict cardinality reduction is possible
exactly when $g_{\max}\ge1$, and then it requires
\begin{equation}
 q\ge\left\lceil S/g_{\max}\right\rceil.\label{eq:fullfactor}
\end{equation}
Indeed, $B\lceil S/q\rceil\gamma<1$ is equivalent to
$\lceil S/q\rceil\le g_{\max}$. At $(d,m,r,B,S)=(4096,14336,128,1,1024)$
this gives $q\ge43$; at $(1536,8960,384,1,1024)$ it gives $q\ge342$.
These are optimistic generation-volume thresholds, not traffic or quality
claims. The causal grouping restriction still applies.

\subsection{Training and routing}
Under dense synchronous data parallelism, communicating gradients for fewer
persistent parameters can reduce communication volume. The generated values
are functions of controller parameters and are not independent optimizer
parameters. This explains a low-rank communication saving regardless of
whether descriptors are materialized during inference.

A runtime router saves work only when it dispatches the selected branch before
that branch's weights and kernels are consumed. Computing both branches and
then selecting with \texttt{torch.where} provides output selection without the
claimed branch-skipping savings. A quality threshold selected on evaluation
data also gives an optimistic policy estimate. A future router evaluation must
use separate training, calibration, and final evaluation data, fixed routing
rules, and whole-call transfer measurements. This paper makes no population
tail-risk guarantee or operational router claim.

\section{Evidence Protocol and Provenance}
The empirical questions are: (Q1) does conditioning improve reconstruction at a
matched budget; (Q2) does fusion improve the same conditional computation;
(Q3) can software rematerialization preserve the computation and gradients;
and (Q4) do quality and efficiency hold simultaneously in a complete model?
Table~\ref{tab:evidence} makes the available evidence explicit.

\begin{table}[t]
\centering\small
\caption{Evidence classes after the October 22, 2026 audit.}
\label{tab:evidence}
\begin{tabular}{p{.25\linewidth}p{.32\linewidth}p{.33\linewidth}}
\toprule
Evidence & Available artifacts & Supported scope \\
\midrule
Qwen FFN quality & Four per-seed summaries and duplicate aggregate record & Arithmetic and within-archive contrasts; original run provenance incomplete \\
Four-arm v1 timing & One same-run summary per geometry & Recorded medians; no local per-iteration uncertainty \\
Multicore v4 & Timing/equivalence records plus failed stability record & Exploratory timings; reliable speedup claim blocked \\
DDP optimizer timing & Three complete JSON run reports & Short kernel-execution timings; numerical health unknown \\
CPU checkpointing & New source, nine cases, hashes and tests & Actual-module saved-storage and gradient verification \\
Algebra checks & New script and 50 seeded fixtures & Numerical consistency of identities; no model-quality evidence \\
\bottomrule
\end{tabular}
\end{table}

\paragraph{Archive integrity.}
The canonical quality family is
\texttt{asi2\_rank512\_control\_20260922T080042Z}. Its two local files have
SHA-256 hashes recorded in the audit. Their per-seed records agree exactly.
Hash matching establishes integrity of those files, not independent proof of
execution. The canonical directory lacks original per-seed execution reports,
trained checkpoints, and complete immutable data manifests. Several alternative
closing-run report fragments cannot be parsed as complete JSON and are excluded.
The revision preserves historical artifacts and records these limitations.

\paragraph{Quality protocol reported by the archive.}
The target is the layer-16 FFN of Qwen2.5-1.5B-Instruct, with $d=1536,m=8960$,
fp32, ranks 384, 401, and 512, 5000 steps, learning rate $2\times10^{-4}$,
batch size one, and 128 evaluation batches. Training seeds are 2026--2029.
Factor seed 4242 and activation seed 1729 are shared. Consequently these are
four training-seed records from shared data preparation, not four independently
disjoint datasets. The record names separate train/test text files and a
near-duplicate filter threshold 0.82; the available archive does not establish
the exact realized splits or checkpoint identity.

The evaluator reports relative MSE as total squared reconstruction error divided
by target squared magnitude for equal-sized evaluation batches. In notation,
$\mathrm{RelMSE}=\norm{\widehat Y-Y}_F^2/\norm{Y}_F^2$ for a nonzero target.
It measures FFN output reconstruction, not language-model perplexity or task
accuracy. All within-rank contrasts use matching seed labels from one family.
The shown intervals are $\overline\Delta\pm t_{.975,3}s_\Delta/\sqrt4$,
with $t_{.975,3}=3.1824463$. They are unadjusted descriptive intervals under an
independent normal paired-run model, whose adequacy cannot be established with
four archived summaries. No preregistered equivalence margin, population
claim, or multiple-comparison significance claim is made.

\section{Controlled Results}
\subsection{Rank, capacity, and conditioning}
Table~\ref{tab:quality} recomputes every value from the same four seed records.
Parameter storage is reconstructed using Eq.~\eqref{eq:controller-count} and
the code-default hidden width $h=192$; this assumption is explicit because the
canonical record omits the resolved width. It is not a measurement of peak
memory. A fixed additive 2 MiB controller approximation is inaccurate across
ranks and is not used.

\begin{table}[t]
\centering\small
\caption{Archived layer reconstruction: mean $\pm$ sample standard deviation
across four training-seed records. Lower MSE is better. Storage includes
factors and, when applicable, the controller under the stated $h=192$
assumption. All rows are recomputed summaries, not newly executed training.}
\label{tab:quality}
\input{\paperroot/generated_tables/journal_quality_rows.tex}
\end{table}

At rank 384 the generated-minus-persistent difference is $-0.0007545$;
its descriptive interval excludes zero (Table~\ref{tab:paired}). This numerical
finding does not establish a meaningful task-quality improvement. The
matched-rank generated-minus-shuffled differences are approximately $-0.000208$
and $-0.000196$ at ranks 384 and 401, about $0.05\%$ of the control MSE.
At rank 512 the difference is approximately $+0.000025$ with an interval spanning zero.
Shuffling is therefore consistent with little observed dependence on the
correct token context in these summaries. It does not prove that the learned
controller is constant, or predict its behavior on other distributions.

\begin{table}[t]
\centering\small
\caption{Paired generated-minus-control MSE contrasts. Negative values favor
the generated arm. Intervals are descriptive and unadjusted; four records do
not certify equivalence, generalization, or independence.}
\label{tab:paired}
\input{\paperroot/generated_tables/journal_contrast_rows.tex}
\end{table}

The closest parameter-budget comparison has generated rank 384 at 48.099 MiB
and persistent rank 401 at 48.167 MiB. Their mean relative MSE values are
0.448878 and 0.440310: the generated arm has a 1.95\% relative penalty.
These are near-matched rather than exactly equal budgets. Increasing persistent
rank to 512 further reduces MSE, but that comparison consumes a larger budget
and cannot isolate conditioning. Figure~\ref{fig:quality} keeps that capacity
cost visible. The observed pattern motivates the constant-core control and
output-subspace diagnostics; neither has yet been measured on these trained
checkpoints.

\begin{figure}[t]
\centering
\includegraphics[width=.86\linewidth]{\paperroot/figs/journal_quality.pdf}
\caption{Recomputed archive MSE versus reconstructed parameter storage.
Error bars are across-seed sample standard deviations, not corpus-level
confidence intervals. Lines join observed rank settings only and do not imply
measured interpolation. Persistent low rank is a competitive capacity control.}
\label{fig:quality}
\end{figure}

\subsection{Fusion versus materialization}
Table~\ref{tab:v1} uses only same-run pairs in the shipped v1 summary.
The probe has $(d,m,r,h)=(4096,14336,128,128)$, $B=1,S=128$, fp16,
$\alpha=0.05$, and seed 1234 on Ascend 910B2. The described protocol uses
30 warmup launches, 200 timed iterations, and five repeats with synchronized
host timing. The dense matrices in this harness are exact products of the shared factors;
this is an execution control, not an independently trained full-rank checkpoint.
The dynamic arms use the same factors and controller, but their outputs need
not match the unmodulated persistent arm. Those original full reports are
designated remote in the summary;
local per-iteration traces are unavailable. Thus the table reports recorded
medians without newly asserted statistical uncertainty or cache-residency claims.

\begin{table}[t]
\centering\small
\caption{Same-run archived v1 timings in milliseconds. The final column is
fused latency divided by materialized latency. The probe includes host-side
caching of tiling/workspace and the fp16 controller-weight image.}
\label{tab:v1}
\begin{tabular}{rrrrrr}
\toprule
Tokens & Dense & Persistent & Materialized & Fused v1 & Fused/materialized \\
\midrule
128 & 0.4611 & 0.2719 & 0.8509 & 0.8234 & 0.96768 \\
4096 & -- & -- & 1.3932 & 2.6821 & 1.92514 \\
\bottomrule
\end{tabular}
\end{table}

At the probe, fusion reduces recorded latency by 3.23\% relative to
materialization, while it remains $3.03\times$ slower than persistent low rank.
At $T=4096$ it is $1.93\times$ slower than materialization. Both comparisons
matter: the first isolates an implementation difference in a generated model;
the second asks whether that model is efficient against the practical static
baseline. No measured quality-matched deployment improvement follows.

The v4 multicore record contains exploratory v1/v4 timing ratios of
1.13, 1.38, 1.70, 2.14, and 3.03 for $T=128,512,1024,2048,4096$.
However, its 1000-call stability test ends with \texttt{pass=false} and a
nonfinite final error. The cause is unresolved; it cannot be assigned to an
environment transient from this record alone. The audit therefore marks the
stability gate \textbf{FAIL} and excludes v4 from validated efficiency claims.
There is also no matched v4-versus-materialized complete-FFN experiment.
Multiplying ratios from different runs does not close this gap.

\subsection{Optimizer-inclusive execution timings}
Table~\ref{tab:training} recomputes three archived four-device run reports at
$d=4096,m=14336,B=2,S=512$, fp16. The timed region includes zeroing gradients,
forward, squared loss, backward, DDP synchronization, and AdamW update.
There are three warmup and five timed training iterations per repeat. The
reported statistic is the maximum across devices of each device's mean time,
then summarized across three repeats. It is not a sample of independent
end-to-end training runs to convergence.

\begin{table}[t]
\centering\small
\caption{Archived optimizer-inclusive synchronized execution, mean $\pm$
sample standard deviation across three reports. Ratios use means. These
microbenchmarks lack recorded finite-loss/gradient/update health gates, so they
support execution-time accounting only.}
\label{tab:training}
\input{\paperroot/generated_tables/journal_timing_rows.tex}
\end{table}

The dense mean is 24.5906 ms, and low-rank ratios range from 3.372 to 4.140.
The benchmark changes parameter count and controller behavior and does not
run the inference descriptor no-store kernel. These ratios cannot identify a
lifetime effect. Moreover, the FP16 optimizer benchmark records no finite-loss,
finite-gradient, or finite-parameter gate. The revision therefore does not
interpret those timings as healthy training, convergence acceleration, or a
quality-preserving replacement of an 8B-parameter language model. The phrase
``8B-scale geometry'' refers only to matrix dimensions.

\subsection{Fresh CPU rematerialization verification}\label{sec:cpu}
To provide complete local evidence for one narrower claim, we executed the
actual \texttt{LowRankFFN} with \texttt{token\_diag\_core}, comparing ordinary
autograd with non-reentrant PyTorch checkpointing. Fixtures use synthetic
random inputs and randomly initialized factors, $d=32,m=64,r=8,h=16$,
$B=2,S=16$, fp32, and $\alpha=0.2$. Residual stacks
$x\mapsto x+\widetilde f(x)$ have depths 1, 2, and 4; seeds are 2026--2028.
Each arm begins from an identical deep copy. Conditioning response and
controller gradients are required to be nonzero.

Saved-tensor hooks count unique underlying nonparameter storages retained by
the FFN forward. Aliases of parameters are excluded; shared storage is counted
once. Caller-owned inputs and full storage underlying saved views are included.
Loss-function saved tensors are outside the measured region. This quantity is
not allocator peak memory, total training memory, or hardware traffic.

\begin{table}[t]
\centering\small
\caption{New CPU software measurements, three seeds per depth. Exact byte
counts were identical across these seeds. Output, loss, input-gradient and
all parameter-gradient maximum absolute differences were zero in all nine
cases. The controller was active and all checked tensors were finite.}
\label{tab:cpu}
\begin{tabular}{rrrrr}
\toprule
Depth & Eager saved bytes & Checkpoint saved bytes & Ratio & Cases \\
\midrule
1 & 53,248 & 4,096 & 13.0 & 3 \\
2 & 106,496 & 8,192 & 13.0 & 3 \\
4 & 212,992 & 16,384 & 13.0 & 3 \\
\bottomrule
\end{tabular}
\end{table}

Python 3.12.14 and PyTorch 2.14.1+cpu were used with one CPU thread. The raw
JSON records exact source hashes, initial-state hashes, per-case tensor checks,
and the measurement boundary. Tests reject missing or nonfinite gradients
and verify parameter-alias exclusion. Table~\ref{tab:cpu} supports software
rematerialization for these fixtures. It supplies no evidence about NPU
checkpoint speed, peak memory, or trained-model quality. The previously quoted
24.5--71.8-fold memory headline is not carried forward because the available
legacy evidence does not establish an instrumented actual-MWG measurement.

\subsection{Algebra checks}
Fifty seeded float64 NumPy fixtures verify the output-subspace identity and
compare the constant-core normal equations with an independently constructed
least-squares system. Maximum rowspace residual is $1.42\times10^{-14}$;
maximum relative orthogonal-decomposition residual is $8.88\times10^{-16}$;
constant-core coefficients agree within $3.67\times10^{-15}$. The fitted
constant has no greater calibration loss than the unweighted mean core in
these fixtures. These are checks of the derivations, explicitly separated
from real-model experiments and from the proofs themselves.

\section{Limitations and Decisive Follow-Up Experiments}
The theoretical analysis is general within its stated architecture and
scheduling assumptions. The available quality evidence is narrow and has
incomplete execution provenance. The accelerator evidence comes from one device
family, with small repeat counts and different configurations for quality and
systems probes. The new CPU experiment establishes a local software property.
These pieces do not establish Q4, a simultaneous full-model quality--latency
benefit.

A decisive next study should freeze a trained causal checkpoint, its tokenizer,
dataset splits, and evaluation harness; then measure both quality and runtime
for that same artifact. The primary quality comparison should include trained
persistent low rank at exact or explicitly near-matched parameter and compute
budgets, a calibrated constant-core control, and relevant quantized baselines.
Held-out perplexity and task metrics must accompany FFN reconstruction. Rank
floors require stored evaluation activations and target outputs; a controller
claim should explain the reducible error above those floors.

For execution, the generated/materialized and generated/fused arms must share
checkpoint, input, precision, warmup, cache policy, and complete-FFN timing
boundaries. Alternate paired launches, report timing distributions, and verify
against an independent high-precision reference before and after extended
execution. Gate all timing claims on finite outputs, gradients where relevant,
healthy repeated calls, and supported shapes/devices. Record device counters
and distinguish logical transfers from DRAM measurements. Test true cached
decode ($S=1$), prefill, varying batch sizes, and the serving stack if a serving
claim is made. The current failed v4 stability check must be resolved before
that variant can enter an efficiency comparison.

Finally, choose quality margins, primary workloads, repeat counts, and stopping
rules before collecting the next comparison. A successful result requires both
acceptable held-out quality and an observed efficiency gain on the same
implementation. If constant-core or persistent-factor controls dominate, the
correct conclusion is a limitation of dynamic modulation in that regime.
The mechanism and negative evidence remain useful, but they should not be
renamed as an established deployment advantage.

\section{Conclusion}
Ephemeral diagonal cores expose a narrow and testable execution choice within
a familiar conditional low-rank model. Their fixed final factor imposes an
output-subspace floor, and eliminating a descriptor tensor removes only the
transfers associated with that tensor under a matched schedule. The archive
shows modest same-rank reconstruction differences, a near-budget disadvantage,
and no validated complete-model efficiency gain. New actual-module CPU checks
establish a reproducible rematerialization property with exact gradient
agreement. These results identify what is supported and which joint experiment
would determine whether the approach is useful in deployment.

\section*{Reproducibility and artifact availability}
The canonical manuscript, audit script, algebra checks, CPU lifecycle test,
and archived source records are in the accompanying research directory.
\texttt{scripts/build\_journal.sh} rebuilds the audit-derived tables, figures,
and author-review PDFs. The audit JSON contains exact input hashes and an
explicit false submission-readiness flag for the unresolved empirical gates.
\texttt{paper/REVIEW\_20261022.md} maps each review finding to its correction or
remaining requirement. Data and checkpoint availability limitations are stated
in the evidence protocol, rather than treating a file hash as execution proof.

\section*{AI assistance disclosure}
Generative AI tools assisted technical review, mathematical derivations,
implementation, verification, analysis, and manuscript revision. Newly executed
synthetic fixtures are labeled as such; archived accelerator records were
reanalyzed without claiming new device execution. The author is responsible for
reviewing the arguments, artifacts, and disclosures before any submission.

\section*{Broader impact}
The intended application is more efficient neural-network execution. Energy
savings, production cost reductions, and improved model capabilities have not
been measured here. Incomplete numerical checks or overstated performance
claims could mislead deployment decisions, which motivates explicit evidence
boundaries and preservation of failed gates.
